\section{Corrections to the original formulation and the retained research programme}
\label{dev:originalcorrection119}

This section records mathematical corrections to the original author source
\texttt{bf049e302e9ed38303084e8da16b13717b052a70}, identified in the
12 major comments of the report dated 15 September 2026.  The original files,
counterexamples, and subsequent numerical records remain in the immutable
development history.  Preservation of an original statement is not an
endorsement of its validity.  In particular, the original composite-loss
optimality proposition, residual-based viscosity-selection assertion, and
unverified equilibrium claims are not premises of any current theorem.
The current Bellman construction, actor comparison, and continuous-risk
certificates are stated and proved separately in the main text and supplement.
The following corrections are elementary diagnostics, not new contributions
or reruns of the original economic experiments.

\subsection{The original scalar loss and the current computational graph}
\label{dev:compositeloss119}

The original loss was
\[
 L(\pi,V)=-\mathbb E H^\pi[V]
       +\omega\mathbb E(-V_t-H^\pi[V])^2.
\]
Joint minimization does not force the second term to vanish.  For the
one-action, no-state problem with terminal value one at date one and
aggregator $f(v)=-v$, the correct solution is $V^*(t)=e^{t-1}$.
With uniform sampling and $\omega=1$, put
$V_\varepsilon(t)=V^*(t)-\varepsilon(1-t)$.
The terminal value is unchanged, whereas
\[
 -V_\varepsilon'+V_\varepsilon=-\varepsilon(2-t),\qquad
 L(V_\varepsilon)-L(V^*)
       =-\frac{\varepsilon}{2}+\frac{7\varepsilon^2}{3}.
\]
This difference is strictly negative for $0<\varepsilon<3/14$ and equals
$-3/112$ at $\varepsilon=3/28$.  Thus the claimed local, and hence global,
optimality assertion was false even with exact terminal data.  In the
stationary scalar case the same issue is the identity
\[
 -H+\omega H^2
   =\omega\left(H-\frac{1}{2\omega}\right)^2
      -\frac{1}{4\omega}.
\]
Its minimizer is not the Bellman root.

The current actor is proposed from a specified fitted continuation and
accepted only using the original objective and the stated independent
certificate.  The current critic fitting criterion has its own parameter
block.  Neither computation is justified by the false joint-loss proposition.
The supplement's computational-graph subsection following
Proposition~\ref{prop:investmentfoundations119} gives the full chain rule for
the criterion $\tfrac12\mathbb E r^2$, including generator and shared-parameter
channels, and states exactly when a detached continuation changes that
derivative.  A zero target or an automatic-differentiation implementation does
not prove optimizer stability, two-timescale convergence, or a utility
gradient under an arbitrary state-sampling law.

\subsection{Boundary data and viscosity selection}
\label{dev:viscosity119}

Zero interior residual does not choose terminal or lateral data: with zero
flow and dynamics, every constant has zero interior residual, while terminal
payoff zero selects only the zero constant.  Proposition~\ref{prop:finitefoundations119}
specifies the terminal function, admissibility, continuous Bellman class,
measurability, and policy recursion actually used in the current finite-horizon
model.  Proposition~\ref{prop:stoppedbsde119} specifies the separate Brownian
boundary/terminal discrepancy under stopping.  An infinite-horizon claim
additionally requires the stated function class and contraction or
transversality hypotheses; it does not follow from a finite-horizon argument.

Even exact boundary data and vanishing population squared residual do not
give a general viscosity-selection theorem.  For the deterministic exit
problem with speed in $[-1,1]$ and payoff minus elapsed time, the equation is
$1-|V'|=0$ on $(-1,1)$ with zero boundary values.  Its value is
$V^*(x)=|x|-1$.  The smooth functions
\[
 W_\varepsilon(x)=\sqrt{1+\varepsilon^2}
                      -\sqrt{x^2+\varepsilon^2}
\]
have the same boundary values and satisfy
\[
 \mathbb E\{1-|W_\varepsilon'|\}^2
 =2-2\sqrt{1+\varepsilon^2}+2\varepsilon
       -\varepsilon\arctan(1/\varepsilon)\longrightarrow0.
\]
Nevertheless, their uniform limit is $W(x)=1-|x|$.  The test function
$\varphi(x)=1$ touches $W$ from above at zero and gives
$1-|\varphi'(0)|=1>0$, violating the viscosity subsolution inequality.
The displayed integral follows by expanding the square on $[0,1]$ and
integrating $x/\sqrt{x^2+\varepsilon^2}$ and
$\varepsilon^2/(x^2+\varepsilon^2)$ separately.

The current results therefore do not identify spectral bias or fixed-network
training with vanishing viscosity.  Uniform action comparisons and coverage
allowances are explicit inputs to the current certificates.  A theorem on
approximating a smooth target and its derivatives does not approximate
nonexistent classical derivatives at a kink, prove optimization of a fixed
architecture, or provide a PDE comparison principle.

\subsection{Analytical portfolio and recursive-utility diagnostics}
\label{dev:mertonez119}

Under the original unconstrained Merton parameters
$(\mu,r,\gamma,\sigma,\rho)=(0.08,0.02,2,0.2,0.04)$, the formal stationary
CRRA benchmark is
\[
 \pi^*=\frac{\mu-r}{\gamma\sigma^2}=\frac34,\qquad
 m=\frac{\rho+(\gamma-1)
       [r+(\mu-r)^2/(2\gamma\sigma^2)]}{\gamma}
       =\frac{33}{800}.
\]
With $\mu=0.01$ the unconstrained allocation is $-1/8$ and its no-short-sale
projection is zero at every positive wealth level.  These values replace the
incorrect analytical numbers in the original table.  A finite horizon needs
its own terminal/bequest specification and generally does not have this
constant consumption ratio.  A control bound gives action-space optimality
conditions; it does not by itself create a wealth threshold, a value-gradient
jump, or an obstacle variational inequality.  The old constant-parameter
example consequently does not validate a switching-boundary claim.

For Epstein--Zin notation set $a=1-1/\psi$, $b=1-\gamma$, and $Y=bV>0$.
The originally printed aggregator was
\[
 f_{\rm old}(c,V)=\frac{\rho b}{a}c^aY^{1-a/b}-\rho V,
 \qquad (f_{\rm old})_c=\rho b c^{a-1}Y^{1-a/b}.
\]
It has negative consumption derivative when $b<0$, including the reported
$\gamma=5$ regime.  Under the conventional difference-form normalization,
the corrected expression is
\[
 f(c,V)=\frac{\rho}{a}
       \{c^a(bV)^{1-a/b}-bV\},\qquad
 f_c=\rho c^{a-1}(bV)^{1-a/b}>0
\]
for $c>0$ and $bV>0$.  The choice of normalization and its domain are part
of the model.  With the original return parameters, $\gamma=5$,
$\psi=3/2$, and $\rho=1/25$, formal stationary substitution gives
\[
 \pi=\frac3{10},\qquad
 m=\psi\rho+(1-\psi)
       \left[r+\frac{(\mu-r)^2}{2\gamma\sigma^2}\right]
       =\frac{91}{2000}.
\]
These algebraic identities are not an infinite-horizon existence or
verification theorem.  They do not establish integrability, admissibility,
selection of a recursive-utility solution, or preservation of the fractional
power domain by an arbitrary network.  The current bounded Lipschitz BSDE
bridge does not assert those properties for the singular Epstein--Zin
specification.  Current monetary-risk Bellman recursions and general nonlinear
recursive utility retain their different hypotheses and objectives.

\subsection{The expected stochastic-Hessian objective}
\label{dev:traceobjective119}

If $\widehat q$ is conditionally unbiased for the diffusion trace $q$, then
\[
 \mathbb E[(a-\widehat q/2)^2\mid a,q]
       =(a-q/2)^2+\tfrac14\operatorname{Var}(\widehat q\mid a,q).
\]
The conditioning here fixes the network, state, and all non-probe
quantities.  The variance usually depends on the network derivatives.
For one standard Gaussian probe with $\widehat q=hZ^2$ and $a=1$, Gaussian
second and fourth moments give $1-h+3h^2/4$, minimized at $h=2/3$, whereas
the true square is minimized at $h=2$.  Two conditionally independent probe
banks give the unbiased product
$\mathbb E[(a-\widehat q_1/2)(a-\widehat q_2/2)]=(a-q/2)^2$.
The individual product can be negative and its variance remains a numerical
cost.  For covariance $\sigma\sigma^\top$, a symmetric matrix actually
probed is $\sigma^\top D^2V\sigma$; for a nonsymmetric representation, the
quadratic form uses its symmetric part.  Unbiasedness of an unweighted trace
does not validate a diffusion-weighted implementation or a squared loss.
No current theorem concludes feasibility failure from trace noise when the
executed actor is hard-mapped into the feasible box.

\subsection{Preference control, temporal selves, and strategic deviations}
\label{dev:extensions119}

For the original additive preference diffusion
$U_t=u_0+\int_0^t\theta_s\,ds+\sigma W_t$, with $\sigma>0$ and
$|\theta_s|\le M$, one has $U_t\le u_0+Mt+\sigma W_t$.
Thus $\mathbb P(U_t<0)>0$ for every $t>0$ by the nonzero Gaussian lower
tail.  Bounded controls do not make the declared positive state invariant.
Moreover,
\[
 \partial_u\frac{c^{1-u}}{1-u}
   =\frac{c^{1-u}\{1-(1-u)\log c\}}{(1-u)^2}>0
       \quad(c>1,\ u>1).
\]
This calculation refutes the original sign argument; it does not determine
the derivative of the endogenous value.  Subtracting a term that varies
with a controlled preference state changes its rewards.  A boundary rule,
cardinal normalization, consumption units, and an observation model are
therefore indispensable to that research application.  Neither its FOC nor
a policy plot proves a comparative-static sign or econometric identification.
The original investment experiment in the main text does not replace the
independent two-dimensional validation still needed for this preference
model.

The original temporal-self Hamiltonian $\beta u(c)+V_x(r x-c)$ has local
FOC $\beta c^{-\gamma}=V_x$; for $\beta>0$, $\gamma>0$, and fixed $V_x>0$,
its solution $c=(\beta/V_x)^{1/\gamma}$ decreases with a lower $\beta$.
For comparison only, a specified one-step present-bias
objective $u(c)+\beta\delta V_{t+1}(F(x,c))$ instead has derivative
$u'(c)+\beta\delta V_{t+1}'(F)F_c$.  Neither algebra is a derivation of a
continuous-time sophisticated equilibrium.  A deviation must include its
effect on the state and continuation.  The original separate supplement's
$\sup_c\{u(c)-u(c^*)\}$ with all dynamics frozen fails that requirement.
The retained nonlinear-deviation result is explicitly conditional on the
specified recursive game and deviation class; slow target-network updates
alone do not establish subgame perfection.

For the Cournot diagnostic $u_i=q_i(1-q_i-q_j)$, the unilateral FOC is
$1-2q_i-q_j=0$ and the symmetric Nash choice is $q_i=1/3$.
Joint profit maximization permits $q_i=1/4$, where the unilateral derivative
is $1/4>0$.  Cross-player derivatives in a global scalar loss therefore
cannot be called unilateral best responses.  Exact evaluation of a policy
profile is also insufficient: exploitability requires a distinct best-response
comparison.  Finite-action certificates retain their finite action set;
continuous-game, equilibrium-existence, and large-market limits require
their own approximation, deviation, and market-scaling arguments.

Finally, the original plotting placeholders remain historical illustrations,
not observations or executed comparison runs.  New source hashes, protocols,
raw results, independent checks, and failures are identified by their own
execution episode.  None of the calculations in this section regenerates the
original portfolio, preference-control, temporal-self, or game experiment.
